\documentclass{article}

\usepackage{amsmath,amssymb}
\usepackage{xurl}
\usepackage[hidelinks]{hyperref}

% Shared commands for notes in docs/paper-gaps/.

\DeclareUnicodeCharacter{2080}{\ifmmode{}_{0}\else\textsubscript{0}\fi}
\DeclareUnicodeCharacter{2081}{\ifmmode{}_{1}\else\textsubscript{1}\fi}
\DeclareUnicodeCharacter{2082}{\ifmmode{}_{2}\else\textsubscript{2}\fi}
\DeclareUnicodeCharacter{2099}{\ifmmode{}_{n}\else\textsubscript{n}\fi}

\newcommand{\Lean}{\textsc{Lean}}
\ExplSyntaxOn
\NewDocumentCommand{\leanid}{m}{
  \ifmmode
    \textsf{\detokenize{#1}}
  \else
    \group_begin:
    \urlstyle{sf}
    \tl_set:Nn \l_tmpa_tl {#1}
    \tl_replace_all:Nnn \l_tmpa_tl {\_} {_}
    \exp_args:NV \path \l_tmpa_tl
    \group_end:
  \fi
}
\ExplSyntaxOff
\providecommand{\tr}{\operatorname{tr}}
\newcommand{\tnleanIssueBase}{https://github.com/LionSR/TNLean/issues/}
\newcommand{\tnleanPrBase}{https://github.com/LionSR/TNLean/pull/}
\newcommand{\tnleanDocBase}{https://sirui-lu.com/TNLean/docs/}
\newcommand{\ghissue}[1]{\href{\tnleanIssueBase#1}{GitHub issue \##1}}
\newcommand{\ghpr}[1]{\href{\tnleanPrBase#1}{PR \##1}}
\newcommand{\doclink}[1]{\href{\tnleanDocBase#1.html}{\path{#1}}}

% Dirac notation and the identity operator, as in blueprint/src/macros/common.tex.
% \providecommand keeps note-local definitions authoritative.
\usepackage{dsfont}
\providecommand{\ket}[1]{|#1\rangle}
\providecommand{\bra}[1]{\langle #1|}
\providecommand{\braket}[2]{\langle #1 | #2 \rangle}
\providecommand{\ketbra}[2]{|#1\rangle\!\langle #2|}
\providecommand{\Id}{\mathds{1}}
\providecommand{\one}{\mathds{1}}
\providecommand{\C}{\mathbb{C}}
\providecommand{\MN}[1]{M_{#1}(\C)}
\providecommand{\MD}{M_D(\C)}

% External referencing for paper sources.  The build helper writes sanitized
% auxiliary files under build/paper-aux/ and excludes duplicated source labels.
\usepackage{xr}

% Import a generated paper auxiliary file with a caller-supplied label prefix.
% The candidate paths support builds from docs/paper-gaps/, the repository root,
% and build/paper-aux/ itself.
\newcommand{\importpaperaux}[2]{%
  \IfFileExists{../../build/paper-aux/#2.aux}{%
    \externaldocument[#1]{../../build/paper-aux/#2}%
  }{%
    \IfFileExists{build/paper-aux/#2.aux}{%
      \externaldocument[#1]{build/paper-aux/#2}%
    }{%
      \IfFileExists{#2.aux}{%
        \externaldocument[#1]{#2}%
      }{}%
    }%
  }%
}

% General external reference macros
\newcommand{\paperref}[2]{\ref{#1-#2}}
\newcommand{\papereqref}[2]{(\ref{#1-#2})}

% CPSV16 source (1606.00608 / MPDO-22-12-17-2.tex) external document and shortcuts
\importpaperaux{cpsv16-}{1606.00608/MPDO-22-12-17-2-tnlean}
\newcommand{\cpsvref}[1]{\paperref{cpsv16}{#1}}
\newcommand{\cpsveqref}[1]{\papereqref{cpsv16}{#1}}

% Verdict marker: \gapnote{<kind>}{<status>}. Kind is the policy
% classification (clarification, local-correction, scope-restriction,
% unfaithful, false-source, open-gap); status is open, wip (elimination
% underway), resolved, or historical. Machine-read by the site index;
% prints a verdict line.
\newcommand{\gapnote}[2]{\par\noindent\textbf{Verdict:} #1 (#2).\par}


\title{The Repeated Case in the Scalar Three-Cocycle Normalization}
\author{}
\date{2026-08-28}

\begin{document}
\maketitle
\gapnote{local-correction}{resolved}

\section{Setting}

Let $G$ be a group with identity element $e$, and write $\C^\times$ for the
multiplicative group of nonzero complex numbers. A scalar $3$-cochain is a
function $\omega\colon G^3\to\C^\times$, written $\omega(g,h,k)$; in
Ref.~\cite{FrancoRubio2025MPUGauging} it is the anomaly $3$-cocycle attached
to the fusion tensors of a matrix product unitary representation of $G$. The
question at issue is which values of $\omega$ the paper's standing gauge
conventions set to $1$.

\section{Printed normalization}

In its list of standing gauge conventions,
Ref.~\cite{FrancoRubio2025MPUGauging} states that every fusion tensor
involving $e$ is trivial and that $\omega$ is consequently normalized. The
accompanying display prints
\begin{align}
    \omega(g,h,e)=\omega(g,h,e)=\omega(g,h,e)
        &=1,
        \qquad \text{for every }g,h.
    \label{eq:fbc25_three_cocycle_normalization_typo_printed}
\end{align}
Thus the same identity-axis case appears three times, while the other two
identity axes do not appear.

\section{Corrected convention}

The preceding sentence requires triviality whenever any fusion input is the
identity. The normalized scalar $3$-cochain convention is therefore
\begin{align}
    \omega(e,h,k)=\omega(g,e,k)=\omega(g,h,e)
        &=1,
        \qquad \text{for every }g,h,k.
    \label{eq:fbc25_three_cocycle_normalization_typo_corrected}
\end{align}
Equation~\ref{eq:fbc25_three_cocycle_normalization_typo_corrected} replaces the
three repeated copies in
Eq.~\ref{eq:fbc25_three_cocycle_normalization_typo_printed} by the three
possible identity-axis cases. It changes neither the scalar codomain nor the
$3$-cocycle equation.

\section{Formal treatment}

The formal definition uses
Eq.~\ref{eq:fbc25_three_cocycle_normalization_typo_corrected}.
This is a local correction of the repeated case in the printed display, not an
additional hypothesis on scalar $3$-cochains.\footnote{Formal declaration:
\leanid{TNLean.Algebra.ScalarThreeCochain.IsNormalized} in
\doclink{TNLean/Algebra/ScalarThreeCocycle}.}

\textbf{Verdict: resolved local fix.}
The source's surrounding standing convention determines all three
identity-axis conditions, so the repetition is typographical.

\bibliographystyle{alpha}
\bibliography{references}

\end{document}